\subsection{序列层滑动块因子：由 \texorpdfstring{$\Fold_m$}{Fold\_m} 诱导的 \texorpdfstring{$\Phi_m$}{Phi\_m}}\label{subsec:Phi_m-sliding-block-factor}

本节处理如下结构性问题：$\Fold_m$ 本身是\emph{窗口层}映射 $\Omega_m\to X_m$，并不直接对应双边子移位之间的因子映射；但将其施加于双边序列的滑动窗口，可得到与时间平移相容（移位等变）的序列层算子，从而可在符号动力系统框架内讨论其像空间的转移结构。

记 $\sigma$ 为左移位：对 $s\in\{0,1\}^{\ZZ}$，定义 $(\sigma s)_t:=s_{t+1}$。将有限集合 $X_m$ 视为字母表，记 $X_m^{\ZZ}$ 的左移位同样记为 $\sigma$。

\begin{definition}[序列层折叠（滑动块码）]\label{def:Phi_m}
对任意分辨率 $m\ge 1$，定义映射
\[
\Phi_m:\{0,1\}^{\ZZ}\to X_m^{\ZZ},\qquad
(\Phi_m(s))_t:=\Fold_m\bigl(s_t s_{t+1}\cdots s_{t+m-1}\bigr).
\]
\end{definition}

\begin{proposition}[记忆/预期性（窗口半径）]\label{prop:Phi_m-radius}
$\Phi_m$ 是一个滑动块码，其\textbf{记忆}为 $0$、\textbf{预期性}为 $m-1$：即对每个时刻 $t$，输出符号 $(\Phi_m(s))_t$ 仅依赖于输入坐标
\(
s_t,s_{t+1},\dots,s_{t+m-1}
\)。
\leanverified{paper\_phi\_m\_radius}
\end{proposition}

\begin{proposition}[移位等变性（时间相容）]\label{prop:Phi_m-equivariant}
对任意 $m\ge 1$ 有 $\Phi_m\circ\sigma=\sigma\circ\Phi_m$。
\leanverified{paper\_fold\_phi\_equivariant\_seeds}
\leanverified{paper\_fold\_phi\_equivariant\_package}
\end{proposition}

\begin{definition}[像空间]\label{def:Y_m}
定义 $\Phi_m$ 的像
\[
Y_m:=\Phi_m\bigl(\{0,1\}^{\ZZ}\bigr)\ \subset\ X_m^{\ZZ}.
\]
\end{definition}

\begin{remark}[不要将 $Y_m$ 与“黄金分割移位”混淆]\label{rem:Ym-not-golden-mean}
需要强调：$X_m$ 是“禁止相邻 $11$”语法在长度 $m$ 上的合法词集合（字母表为 $X_m$），而 $Y_m$ 是字母表 $X_m$ 上的一个子移位；因此 $Y_m$ 并不是字母表 $\{0,1\}$ 上的黄金分割 SFT（禁词 $11$ 的 golden mean shift）。

更具体地：
\begin{itemize}
  \item 当 $m=1$ 时，$X_1=\{0,1\}$ 且 $\Fold_1$ 为恒等，因此 $\Phi_1$ 是 $\{0,1\}^{\ZZ}$ 上的恒等映射，从而 $Y_1=\{0,1\}^{\ZZ}$（全移位），这与黄金分割移位不同。
  \item 当 $m\ge 2$ 时，字母表大小为 $|X_m|=F_{m+2}>2$，故 $Y_m$ 的自然字母表不是 $\{0,1\}$；其结构需在 $X_m^{\ZZ}$ 上讨论，并由 sofic 图像理论刻画（见命题 \ref{prop:Y_m-sofic} 与定理 \ref{thm:Phi_m-sofic-graph}）。
\end{itemize}
\end{remark}

\begin{proposition}[对子移位的因子性质（序列级折叠的最小口径）]\label{prop:Phi_m-factor-on-subshift}
设 $S\subset\{0,1\}^{\ZZ}$ 为任意子移位（紧致、移位不变闭集）。则：
\begin{enumerate}
  \item $\Phi_m|_S:S\to X_m^{\ZZ}$ 连续且与移位交换；
  \item 其像 $Y_m(S):=\Phi_m(S)$ 是 $X_m^{\ZZ}$ 的一个子移位；
  \item 映射 $\Phi_m|_S:S\twoheadrightarrow Y_m(S)$ 是一个因子映射。
\end{enumerate}
\leanverified{ShiftInvariant}
\leanverified{shiftInvariant\_univ}
\leanverified{shiftInvariant\_empty}
\leanverified{slideBlockCode\_restriction\_equivariant}
\leanverified{slideBlockCode\_image\_shiftInvariant}
\leanverified{slideBlockCode\_restrict\_surjective}
\leanverified{paper\_Phi\_m\_factor\_on\_subshift}
\end{proposition}

\begin{proposition}[像空间为 sofic 子移位]\label{prop:Y_m-sofic}
$Y_m$ 是一个 sofic 子移位。更具体地，它可由一个显式有限标记图给出。
\end{proposition}

% --- Distillation writeback ---
\begin{proposition}[\texorpdfstring{$\Phi_m$}{Phi m} 像空间的有限区间 incidence 账本]\label{prop:distill-euclid-elements-phi-finite-interval-ledger-compactness}
\textbf{Dependency status: closure.}
固定 \(m\ge1\)。对任意 \(y\in X_m^{\ZZ}\) 与任意有限整数区间
\[
I=[r,s]\cap\ZZ
\qquad(r\le s),
\]
定义有限区间约束集
\[
K_I(y):=\left\{u\in\{0,1\}^{\ZZ}:
\Fold_m(u_tu_{t+1}\cdots u_{t+m-1})=y_t\ \text{for every }t\in I
\right\}.
\]
则
\[
y\in Y_m
\quad\Longleftrightarrow\quad
K_I(y)\ne\varnothing\ \text{for every finite interval }I\subset\ZZ.
\]
等价地，\(Y_m\) 的成员资格可由所有有限区间上的 typed incidence 账本共同闭合；有限张图的相容性通过区间外包而非任意有限并来核验。
\end{proposition}

% --- Distillation writeback ---
% Distillation timestamp: 2026-04-23T12:04:19+08:00; pipeline index: 36.
\begin{proposition}[\texorpdfstring{$\Phi_m$}{Phi m} 局部窗口 lift 的重叠审计]\label{distill:euclid_elements-lifted_congruence_and_residual_rigidity_families-phi-local-window-overlap-lift-audit}
\textbf{Dependency status: closure.}
固定 \(m\ge1\)。设
\[
J=[a,b]\cap\ZZ
\]
为有限整数区间，并令 \(\{I_\alpha\}_{\alpha\in A}\) 为 \(J\) 的有限区间覆盖。对区间
\(I=[r,s]\cap\ZZ\)，记其 \(m\)-窗口支撑膨胀为
\[
I^{+}:=[r,s+m-1]\cap\ZZ.
\]
给定 \(y\in X_m^{\ZZ}\)。称一族有限词
\[
u_\alpha\in\{0,1\}^{I_\alpha^{+}}\qquad(\alpha\in A)
\]
为 \(y\) 在覆盖 \(\{I_\alpha\}\) 上的局部 lift 账本，若对每个 \(\alpha\) 与每个
\(t\in I_\alpha\) 都有
\[
\Fold_m\bigl(\nu_\alpha(t)\nu_\alpha(t+1)\cdots\nu_\alpha(t+m-1)\bigr)=y_t.
\]
称该账本通过重叠 lift 审计，若对任意 \(\alpha,\beta\) 都有
\[
\nu_\alpha|_{I_\alpha^{+}\cap I_\beta^{+}}
=
\nu_\beta|_{I_\alpha^{+}\cap I_\beta^{+}}.
\]
则
\[
K_J(y)\ne\varnothing
\]
当且仅当存在一份通过重叠 lift 审计的局部 lift 账本。特别地，所有单片 \(K_{I_\alpha}(y)\) 非空并不推出 \(K_J(y)\ne\varnothing\)；唯一缺口正是膨胀重叠上的 lift residual。
\end{proposition}
\begin{proof}
若 \(K_J(y)\ne\varnothing\)，取 \(u\in K_J(y)\)。对每个 \(\alpha\)，令
\[
\nu_\alpha:=u|_{I_\alpha^{+}}.
\]
由于 \(u\in K_J(y)\)，对所有 \(t\in I_\alpha\subseteq J\) 都有
\[
\Fold_m(u_tu_{t+1}\cdots u_{t+m-1})=y_t,
\]
故 \(\nu_\alpha\) 是局部 lift。它们来自同一个全局序列 \(u\)，所以在所有膨胀重叠上自动一致，因而通过重叠 lift 审计。

反过来，设存在通过重叠 lift 审计的局部账本 \(\{\nu_\alpha\}_{\alpha\in A}\)。由于它们在所有二重膨胀重叠上一致，可在有限并
\[
J^{+}:=[a,b+m-1]\cap\ZZ=\bigcup_{\alpha\in A} I_\alpha^{+}
\]
上粘合为唯一有限词 \(\nu\in\{0,1\}^{J^{+}}\)。对每个 \(t\in J\)，取 \(\alpha\) 使 \(t\in I_\alpha\)。则窗口
\([t,t+m-1]\cap\ZZ\) 包含在 \(I_\alpha^{+}\) 中，且 \(\nu\) 在该窗口上等于 \(\nu_\alpha\)，故
\[
\Fold_m\bigl(\nu(t)\nu(t+1)\cdots\nu(t+m-1)\bigr)=y_t.
\]
任意把 \(\nu\) 延拓为双边序列 \(u\in\{0,1\}^{\ZZ}\)，便有 \(u\in K_J(y)\)。

最后，若每个单片都有局部 lift，但任意选择的局部 lift 族都在某个膨胀重叠上不一致，则上述粘合步骤无法执行；这不是单个窗口中的构造失败，而是局部 lift 在重叠 incidence ledger 上的残余失败。证毕。
\end{proof}

% --- End distillation writeback ---

\begin{proof}
若 \(y\in Y_m\)，则存在 \(u\in\{0,1\}^{\ZZ}\) 使 \(y=\Phi_m(u)\)。于是同一个 \(u\) 属于每个 \(K_I(y)\)，故所有 \(K_I(y)\) 非空。

反过来，假设每个有限区间 \(I\) 上 \(K_I(y)\ne\varnothing\)。每个 \(K_I(y)\) 是有限多个 cylinder 条件的交，因此是紧空间 \(\{0,1\}^{\ZZ}\) 中的闭集。取有限多个区间 \(I_1,\ldots,I_N\)，令
\[
J:=[\min_
u \inf I_\nu,\ \max_
u \sup I_\nu]\cap\ZZ
\]
为包含它们的有限区间外包。由定义有
\[
K_J(y)\subseteq\bigcap_{
u=1}^N K_{I_\nu}(y),
\]
而 \(K_J(y)\ne\varnothing\)。故族 \(\{K_I(y)\}_I\) 具有有限交性质。紧性给出
\[
\bigcap_I K_I(y)\ne\varnothing.
\]
取 \(u\) 属于该交。对每个 \(t\in\ZZ\)，应用单点区间 \(I=\{t\}\)，得到
\[
(\Phi_m(u))_t=\Fold_m(u_t\cdots u_{t+m-1})=y_t.
\]
故 \(\Phi_m(u)=y\)，即 \(y\in Y_m\)。证毕。
\end{proof}

% --- End distillation writeback ---

\leanverified{paper\_Ym\_sofic}
像空间为 sofic 的结论可参见 \cite{LindMarcus1995}；下一个命题给出可计算的有限图构造。
